% !TEX TS-program = lualatex
%\documentclass[handout]{beamer}
\documentclass{beamer}
\input{../common/misomva}
\newcommand{\misoclasstinytitleslide}{Partially based on material by:  Ulrike von Luxburg, \\[-1em] Gary Miller, Doyle \& Schnell, Daniel Spielman}
\newcommand{\misoclassdate}{}
\usetikzlibrary{graphs}

\begin{document}
% Topic-specific title slide

\newcommand{\topictitle}{Spectral Clustering: Relaxation}

\newcommand{\topicsubtitle}{From Discrete to Continuous}

\MakeTopicSlide

%%%%%%%%%%%%%%%%%%%%%%%%%%%%%%%%%%%%%%%%%%%%%%%%%%%%%%%%%%%%

%%%%%%%%%%%%%%%%%%%%%%%%%%%%%%%%%%%%%%%%%%%%%%%%%%%%%%%%%%%%

\begin{frame}
\frametitle{Spectral Clustering: Cuts on graphs}

\begin{center}
% Manual layout for spectral clustering graph (spring layout needs lualatex)
\begin{tikzpicture}[
  nd/.style={circle, draw=mvaAccent, minimum size=6mm, inner sep=0pt, font=\small},
  ed/.style={draw=mvaAccent, fill=mvaAccent, thick}]
\node[nd] (A) at (0,2) {A};
\node[nd] (B) at (1.5,2) {B};
\node[nd] (C) at (0.75,0.8) {C};
\node[nd] (D) at (2.5,0.8) {D};
\node[nd] (E) at (1.5,-0.2) {E};
\node[nd] (F) at (3.5,-0.5) {F};
\node[nd] (G) at (4.5,0.8) {G};
\node[nd] (H) at (4.5,-1.5) {H};
\node[nd] (I) at (2.5,-1.8) {I};
\node[nd] (J) at (3.5,-2.5) {J};
\node[nd] (K) at (3.5,-3.8) {K};
\draw[ed] (A)--(B); \draw[ed] (B)--(C); \draw[ed] (C)--(A);
\draw[ed] (B)--(E); \draw[ed] (E)--(D); \draw[ed] (D)--(C); \draw[ed] (C)--(E);
\only<1| handout:1>{\draw[ed] (E)--(F); \draw[ed] (D)--(G);}
\only<2| handout:2>{\draw[mvaHighlight, thick] (E)--(F); \draw[mvaHighlight, thick] (D)--(G);}
\only<3-| handout:3->{\draw[ed] (E)--(F); \draw[ed] (D)--(G);}
\draw[ed] (F)--(G); \draw[ed] (F)--(I); \draw[ed] (I)--(J); \draw[ed] (J)--(H); \draw[ed] (H)--(F);
\draw[ed] (G)--(H);
\only<1-3| handout:1-2>{\draw[ed] (J)--(K);}
\only<4| handout:3>{\draw[mvaHighlight, thick] (J)--(K);}
\end{tikzpicture}

\end{center}

\begin{flushright}
\only<2 | handout:2>{
\yellbox {Defining the cut objective we get the clustering!}
}

\only<3,4| handout:3>{
\textbf{MinCut}: ${\rm cut}(A,B)= \sum_{i\in A, j\in B} w_{ij} \qquad$ 
\questionbox{Are we done?}  \\
}
\smallskip
\only<4| handout:3>{
{\tiny Can be solved efficiently, but maybe not what we want $\dots$.}
}

\end{flushright}

\end{frame}

%%%%%%%%%%%%%%%%%%%%%%%%%%%%%%%%%%%%%%%%%%%%%%%%%%%%%%%%%%%%

\begin{frame}
\frametitle{Spectral Clustering: Balanced Cuts}
\vspace{-1.5em}
\begin{flushright}
\yellbox{Let's balance the cuts!}
\end{flushright}
\vspace{-1em}
\pause
\begin{block}{MinCut}
\[{\rm cut}(A,B)= \sum_{i\in A, j\in B} w_{ij}\]
\end{block}

\pause

\begin{block}{RatioCut}
\[{\rm RatioCut}(A,B)= \sum_{i\in A, j\in B} w_{ij} \left( \frac{1}{|A|} + \frac{1}{|B|} \right) \]
\end{block}

\pause


\begin{block}{Normalized Cut}
\[{\rm NCut}(A,B)= \sum_{i\in A, j\in B} w_{ij} \left( \frac{1}{{\rm vol}(A)} + \frac{1}{{\rm vol}(B)} \right) \]
\end{block}

\end{frame}
%%%%%%%%%%%%%%%%%%%%%%%%%%%%%%%%%%%%%%%%%%%%%%%%%%%%%%%%%%%%
\begin{frame}
\frametitle{Spectral Clustering: Balanced Cuts}

\[{\rm RatioCut}(A,B)=  {\rm cut}(A,B) \left( \frac{1}{|A|} + \frac{1}{|B|} \right) \]
\[{\rm NCut}(A,B)= {\rm cut}(A,B) \left( \frac{1}{{\rm vol}(A)} + \frac{1}{{\rm vol}(B)} \right) \]
\bigskip
Easily generalizable to $k \geq 2$
\pause
\medskip
\begin{center}
\questionbox{Can we compute this?}
 \pause 
RatioCut and NCut are NP hard :( 
\pause 
\bigskip
\bigskip
\yellbox{Approximate!}
\end{center}

\end{frame}
%%%%%%%%%%%%%%%%%%%%%%%%%%%%%%%%%%%%%%%%%%%%%%%%%%%%%%%%%%%%
%%%%%%%%%%%%%%%%%%%%%%%%%%%%%%%%%%%%%%%%%%%%%%%%%%%%%%%%%%%%

\begin{frame}
\frametitle{Spectral Clustering: Relaxing Balanced Cuts}
\begin{block}{Relaxation for (simple) balanced cuts for 2 sets}
\[\min_{A,B} {\rm cut}(A,B) \quad {\rm s.t.} \quad  |A|=|B|\]
\end{block}
\bigskip \pause
Graph function $\bff$ for cluster membership\pause : 
$ f_i = 
\begin{cases}
1 & \mbox{if } V_i \in A, \\
-1 & \mbox{if } V_i \in B. \\ 
\end{cases}$
\bigskip \pause
\questionbox{What it is the cut value with this definition?}
\[
{\rm cut}(A,B) 
= \pause \sum_{i\in A, j\in B} w_{i,j}
= \tfrac14 \sum_{i,j}  w_{i,j} (f_i-f_j)^2
= \tfrac12 \bff\transpose \bL \bff
\]
\pause
\questionbox{What is the relationship with the \gcol{smoothness} of a graph function?}

\end{frame}
%%%%%%%%%%%%%%%%%%%%%%%%%%%%%%%%%%%%%%%%%%%%%%%%%%%%%%%%%%%%
\begin{frame}
\frametitle{Spectral Clustering: Relaxing Balanced Cuts}
\[
{\rm cut}(A,B) 
= \sum_{i\in A, j\in B} w_{i,j}
= \tfrac14 \sum_{i,j}  w_{i,j} (f_i-f_j)^2
= \tfrac12 \bff\transpose \bL \bff
\]
\pause
$|A| = |B| \implies \sum_i f_i = \pause 0 \implies \pause \bff \perp \bOne_\nodes$ 
\\ \bigskip 
$\|\bff\| = \pause \sqrt{\nodes}$ \pause

\begin{center}
\begin{block}{\centering objective function of spectral clustering}
\[
\min_{\bff} \bff\transpose \bL \bff \quad {\rm s.t.} 
\quad f_i = \pm 1,
\quad \bff \perp \bOne_\nodes, 
\quad \|\bff\| =  \sqrt{\nodes}
\]
\end{block}
\bigskip
 \pause 
Still NP hard :( \pause  $\quad \to \quad$
\yellbox{Relax even further!} \\
\pause \bigskip 
\tikz[baseline]
\node [cross out,draw=red,anchor=text] {$f_i = \pm 1$};
$\quad \to \quad f_i \in \R$

\end{center}
\end{frame}
%%%%%%%%%%%%%%%%%%%%%%%%%%%%%%%%%%%%%%%%%%%%%%%%%%%%%%%%%%%%
\begin{frame}
\frametitle{Spectral Clustering: Relaxing Balanced Cuts}
\begin{block}{\centering objective function of spectral clustering}
\[
\min_{\bff} \bff\transpose \bL \bff \quad {\rm s.t.} 
\quad f_i \in \R,
\quad \bff \perp \bOne_\nodes, 
\quad \|\bff\| =  \sqrt{\nodes}
\]
\end{block}


\only<2,3 | handout:1>{
\begin{block}{\centering Rayleigh-Ritz theorem}
If $\lambda_1 \leq \dots \leq \lambda_\nodes$ are the eigenvalues of real symmetric $\bL$ then
\begin{align*}
 \lambda_1 &= \min_{\bx \ne 0} \frac{\bx\transpose\bL \bx}{\bx\transpose\bx} = \min_{\bx\transpose\bx=1}\bx\transpose\bL \bx\\
 \lambda_\nodes &= \max_{\bx \ne 0} \frac{\bx\transpose\bL \bx}{\bx\transpose\bx} = \max_{\bx\transpose\bx=1}\bx\transpose\bL \bx 
\end{align*}
\end{block}
$\frac{\bx\transpose\bL \bx}{\bx\transpose\bx} \equiv$ \gcol{Rayleigh quotient}
}

\only<3| handout:1>{ \begin{center}\questionbox{How can we use it?} \end{center} }
 
 \only<4| handout:2>{
 \begin{block}{\centering Generalized Rayleigh-Ritz theorem (Courant-Fischer-Weyl)}
 If $\lambda_1 \leq \dots \leq \lambda_\nodes$ are the eigenvalues of real symmetric $\bL$ 
 and $\bv_1, \dots, \bv_\nodes$ the corresponding orthogonal eigenvectors, then for $k = 1:\nodes-1$
 \begin{align*}
  \lambda_{k+1} &= \min_{\bx \ne 0, \bx \perp \bv_1, \dots \bv_k} \frac{\bx\transpose\bL \bx}{\bx\transpose\bx} 
      = \min_{\bx\transpose\bx=1, \bx \perp \bv_1, \dots \bv_k}\bx\transpose\bL \bx\\
  \lambda_{\nodes-k} &= \max_{\bx \ne 0, \bx \perp \bv_n, \dots \bv_{\nodes-k+1}} \frac{\bx\transpose\bL \bx}{\bx\transpose\bx} 
      = \max_{\bx\transpose\bx=1, \bx \perp \bv_\nodes, \dots \bv_{\nodes-k+1}}\bx\transpose\bL \bx 
 \end{align*}
 \end{block}
 }

\end{frame}


%%%%%%%%%%%%%%%%%%%%%%%%%%%%%%%%%%%%%%%%%%%%%%%%%%%%%%%%%%%%



\MakeLastSlide
\end{document}
